\section{第5组：狭义相对论}
\sourceproblem{5-T1}{17}{21}
\problemtext 不稳定粒子的固有寿命为 $\tau_0$，实验室速度为 $u$，求平均寿命对应的飞行距离。扫描页寿命指数较模糊，本条以符号形式保留。
\workings 实验室寿命 $\tau=\gamma\tau_0$，$\gamma=(1-u^2/c^2)^{-1/2}$，故
\[
L=u\gamma\tau_0.
\]
若原页速度读作 $2\times10^8\,\mathrm{m/s}$，取 $c=3\times10^8\,\mathrm{m/s}$，则 $L=(6/\sqrt5)\times10^8\tau_0$（SI）。
\crosscheck 固有时间为粒子自身时钟的时间；原图寿命指数待高分辨率二校，暂不填写未经确认的数值答案。本条标为条件式待核。
\sourceproblem{5-T2}{17}{21}
\problemtext 静止正方形面积100平方米，观察系沿其一条对角线以 $0.8c$ 运动，求观察系同时测量的面积。
\workings 沿运动方向尺寸缩短为 $\sqrt{1-0.8^2}=0.6$ 倍，垂直方向不变；平面面积的线性变换行列式为0.6，故面积为 $60\,\mathrm{m^2}$。
\crosscheck 形状一般变为菱形；不能将两条边均直接乘0.6，也不是照片的视觉面积。
\sourceproblem{5-T3}{18}{22}
\problemtext 刚性尺静止于 $S'$，与 $x'$ 轴成30度；在沿正 $x$ 轴相对运动的 $S$ 中测得与 $x$ 轴成45度，求相对速度。
\workings 同时测量有 $L_x=L'_x/\gamma,L_y=L'_y$，故 $\tan45^\circ=\gamma\tan30^\circ$，$\gamma=\sqrt3$。因此 $u=c\sqrt{2/3}\approx0.816c$。
\crosscheck 运动方向长度收缩使倾角增大；测量两端须在 $S$ 中同时进行。
\sourceproblem{5-T4}{18}{22}
\problemtext 两飞船固有长度各100米；甲船测量自身前端通过乙船全长的时间 $\Delta t$，求相对速率。扫描时标读作 $(8/3)\times10^{-7}\,\mathrm s$，指数保留二校标记。
\workings 甲系内乙船长度为 $L_0/\gamma$，故 $\Delta t=L_0/(\gamma u)$。解得
\[
u=\frac{c}{\sqrt{1+(c\Delta t/L_0)^2}}.
\]
按上述时标，$c\Delta t/L_0=0.8$，故 $u=5c/\sqrt{41}\approx0.781c$。
\crosscheck 两通过事件发生在甲船前端同一位置；不应把甲船全长也加进通过距离。若时标指数复核改变，使用保留的通式重算。
\sourceproblem{5-T5}{19}{23}
\problemtext 粒子在 $S'$ 中速率 $c/2$、与 $x'$ 轴夹60度，初始通过公共原点；$S'$ 沿 $S$ 正 $x$ 轴以 $0.6c$ 运动，求 $S$ 中运动方程。
\workings $u'_x=c/4,u'_y=\sqrt3c/4$，相对论速度变换给出
\[
u_x=\frac{u'_x+0.6c}{1+0.6u'_x/c}=\frac{17c}{23},\qquad
u_y=\frac{u'_y\sqrt{1-0.6^2}}{1+0.6u'_x/c}=\frac{4\sqrt3c}{23}.
\]
故 $x=17ct/23,\ y=4\sqrt3ct/23$。
\crosscheck 合速率 $c\sqrt{337}/23<c$，并可由洛伦兹坐标变换消去 $t'$ 验证。
\sourceproblem{5-T6}{19}{23}
\problemtext 列车速度 $0.6c$，车上A、B沿车轴相距10米且B在前；地面系记录A处事件后12.5纳秒发生B处事件。求车内系事件先后。
\workings 地面同时测得车内间距 $10/\gamma=8\,\mathrm m$，两事件空间差 $\Delta x=8+0.6c\Delta t=10.25\,\mathrm m$。洛伦兹变换
\[
\Delta t'=\gamma(\Delta t-0.6\Delta x/c)=-10\,\mathrm{ns}.
\]
车内系中B处事件早于A处事件10纳秒。
\crosscheck $c\Delta t=3.75\,\mathrm m<\Delta x$，两事件为空间样间隔，先后可随参考系改变；不能把事件先后等同于因果关系。

\sourceproblem{5-T7}{20}{24}
\problemtext 实验室中两个电子沿相反方向以相同速率 $0.67c$ 运动，求相对速率。
\workings 共线速度变换给出 $u_{\rm rel}=2(0.67)c/(1+0.67^2)\approx0.925c$。
\crosscheck 结果小于光速；不能用经典速度相加得到 $1.34c$。
\sourceproblem{5-T8}{20}{24}
\problemtext 原子核以 $0.5c$ 远离观察者，沿运动方向发射相对核速率 $0.8c$ 的电子，又向观察者方向发射光子，求观察者测得两者速率。
\workings 远离方向为正，$u_e=(0.8+0.5)c/(1+0.8\cdot0.5)=13c/14$。光子在核系速度为 $-c$，变换后 $u_\gamma=(-1+0.5)c/(1-0.5)=-c$，速率为 $c$。
\crosscheck 电子速率小于光速，光子速率在惯性系间不变。
\sourceproblem{5-T9}{21}{25}
\problemtext S系光子沿与正 $x$ 轴成60度方向传播，S'沿正 $x$ 以 $0.6c$ 运动，求S'中的方向及速率。
\workings
\[
u'_x=\frac{c/2-0.6c}{1-0.6/2}=-c/7,\quad
u'_y=\frac{(\sqrt3c/2)\sqrt{1-0.6^2}}{1-0.6/2}=4\sqrt3c/7.
\]
所以 $\theta'=\arccos(-1/7)\approx98.21^\circ$，速率为 $c$。
\crosscheck 两分量平方和为 $c^2$；方向在第二象限，不能直接取反正切主值。
\sourceproblem{5-T10}{21}{25}
\problemtext 质量 $m_0$ 的粒子分别从静止加速至 $0.100c$、从 $0.900c$ 加速至 $0.980c$，求外力净功。
\workings 相对论动能 $K=(\gamma-1)m_0c^2$，故
\[
W_1=(1/\sqrt{0.99}-1)m_0c^2\approx0.005038m_0c^2,
\quad W_2=(1/\sqrt{1-0.98^2}-1/\sqrt{1-0.90^2})m_0c^2
\approx2.7307m_0c^2.
\]
\crosscheck 两功均为正；高速区不能使用经典动能差。题面未给静质量数值，因此以 $m_0c^2$ 表示。
\sourceproblem{5-T11}{22}{26}
\problemtext 电子动能1 MeV、静能0.511 MeV，求动量大小。
\workings $E=K+m_0c^2=1.511\,\mathrm{MeV}$，由 $E^2=p^2c^2+m_0^2c^4$ 得
\[
p=\sqrt{1.511^2-0.511^2}\,\mathrm{MeV}/c
=\sqrt{2.022}\,\mathrm{MeV}/c\approx1.422\,\mathrm{MeV}/c.
\]
换算约 $7.60\times10^{-22}\,\mathrm{kg\,m/s}$。
\crosscheck $pc=\sqrt{K(K+2m_0c^2)}$，不能把动能误当总能量。
\sourceproblem{5-T12}{22}{26}
\problemtext 静止质量42u的粒子衰变为两碎片，其中一片静质量20u、速率 $3c/5$，求另一片动量、总能量和静质量。
\workings 第一片 $\gamma=5/4$，$E_1=25\,\mathrm u c^2$、$p_1=15\,\mathrm u c$。由四动量守恒，第二片 $\bm p_2=-\bm p_1$，$E_2=(42-25)\,\mathrm u c^2=17\,\mathrm u c^2$，
\[
m_2=\sqrt{17^2-15^2}\,\mathrm u=8\,\mathrm u .
\]
按题给 $1\,\mathrm u=1.66\times10^{-27}\,\mathrm{kg}$、$c=3\times10^8\,\mathrm{m/s}$，$p_2=7.47\times10^{-18}\,\mathrm{kg\,m/s}$，$E_2\approx2.54\times10^{-9}\,\mathrm J$，$m_2=1.328\times10^{-26}\,\mathrm{kg}$。
\crosscheck 第二片速率 $p_2c^2/E_2=15c/17<c$，两片动量相反，总能量等于母粒子静能。
